% 分野別類題: 2階線形偏微分方程式の分類と標準形 解答例
% コンパイル: TEXINPUTS="../../../00_common:$TEXINPUTS" latexmk -xelatex answers.tex
\documentclass[a4paper,11pt]{article}
\usepackage{preamble}
\usepackage{shoakubox}

\begin{document}

\kamokumei{類題演習 解答例 --- 2階線形偏微分方程式の分類と標準形}

\daimon{【解法】2階線形偏微分方程式の分類と標準形化の解き方と導出}

\noindent
2階線形偏微分方程式（主要部のみ考える）
\[
A(x,y)\,u_{xx} + 2B(x,y)\,u_{xy} + C(x,y)\,u_{yy} + (\text{低階の項}) = 0
\]
の型は、判別式 $D = B^{2} - AC$ の符号によって次のように分類される。
\[
D = B^{2}-AC
\begin{cases}
< 0 & \text{楕円型 (elliptic)}\\
= 0 & \text{放物型 (parabolic)}\\
> 0 & \text{双曲型 (hyperbolic)}
\end{cases}
\]

\medskip
\noindent
\textbf{特性曲線の導出。} 標準形を得るため、新しい変数 $\xi = \xi(x,y)$, $\eta = \eta(x,y)$ への変換を考える。連鎖律により
\[
u_x = u_\xi \xi_x + u_\eta \eta_x, \qquad
u_y = u_\xi \xi_y + u_\eta \eta_y
\]
などとなり、これを $Au_{xx}+2Bu_{xy}+Cu_{yy}$ に代入すると、$u_{\xi\xi}$ の係数は $A\xi_x^2+2B\xi_x\xi_y+C\xi_y^2$ となる。もし $\varphi(x,y)=\mathrm{const.}$ が常微分方程式
\[
A\left(\frac{dy}{dx}\right)^{2} - 2B\frac{dy}{dx} + C = 0
\]
（これを特性方程式という。$\varphi_x + \varphi_y \frac{dy}{dx}=0$、すなわち $\frac{dy}{dx} = -\varphi_x/\varphi_y$ を上式の主要部係数がゼロになる条件 $A\varphi_x^2+2B\varphi_x\varphi_y+C\varphi_y^2=0$ に代入して得られる）の解（特性曲線）であれば、$\xi=\varphi(x,y)$ ととると $u_{\xi\xi}$ の係数が消える。特性方程式を $\frac{dy}{dx}$ について解くと
\[
\frac{dy}{dx} = \frac{B \pm \sqrt{B^{2}-AC}}{A}
\]
となる。
\begin{itemize}
\item \textbf{双曲型（$D>0$）}: 実数の特性曲線が2本 $\varphi_1(x,y)=c_1,\ \varphi_2(x,y)=c_2$ 得られる。$\xi=\varphi_1,\ \eta=\varphi_2$ とおくと $u_{\xi\xi}, u_{\eta\eta}$ の係数がともに消え、標準形 $u_{\xi\eta}=0$（またはこれに低階の項が加わった形）が得られる。$u_{\xi\eta}=0$ の一般解は $u=f(\xi)+g(\eta)$（$f,g$ は任意関数）。
\item \textbf{放物型（$D=0$）}: 特性曲線は1本のみ。$\xi=\varphi_1(x,y)$、$\eta$ は $\xi,\eta$ が独立になるように任意に選ぶ（例えば $\eta=y$ など）。このとき $u_{\xi\xi}$ の係数も $u_{\xi\eta}$ の係数も消え、標準形 $u_{\eta\eta}=0$ が得られる。一般解は $u=f(\xi)+\eta\,g(\xi)$。
\item \textbf{楕円型（$D<0$）}: 特性曲線は複素共役な対 $\xi=\alpha+i\beta,\ \eta=\alpha-i\beta$ となる。$\alpha=\frac{\xi+\eta}{2},\ \beta=\frac{\xi-\eta}{2i}$ を新変数にとると標準形 $u_{\alpha\alpha}+u_{\beta\beta}=0$（ラプラス型）が得られる。
\end{itemize}

\medskip
\noindent
係数が変数の場合、$D=B^2-AC$ の符号が点ごとに変わりうるため、領域を分けて型を判定する必要がある。

\clearpage
\begin{mondaiboxS}{問1}
$y\,u_{xx} + u_{yy} = 0$ の型を判定せよ（領域によって型が変わる場合はその境界も述べよ）。
\end{mondaiboxS}
\kaitou
$A=y,\ B=0,\ C=1$ であるから
\[
D = B^{2}-AC = 0 - y\cdot 1 = -y
\]
よって
\[
D = -y
\begin{cases}
<0 & (y>0) \Rightarrow \text{楕円型}\\
=0 & (y=0) \Rightarrow \text{放物型（退化）}\\
>0 & (y<0) \Rightarrow \text{双曲型}
\end{cases}
\]
（これは Tricomi 方程式と呼ばれる有名な混合型方程式であり、$y=0$ が型の変わる境界である。）
\[
\boxed{\; y>0:\text{楕円型},\quad y=0:\text{放物型},\quad y<0:\text{双曲型}\;}
\]
（検算: $y=1>0$ では $D=-1<0$ で楕円型、$y=-1<0$ では $D=1>0$ で双曲型となり、符号の対応が一致する。）

\clearpage
\begin{mondaiboxS}{問2}
$u_{xx} - 4u_{xy} + 3u_{yy} = 0$ を特性曲線を用いて標準形に変換し、一般解を求めよ。
\end{mondaiboxS}
\kaitou
$A=1,\ B=-2,\ C=3$ であるから
\[
D = B^{2}-AC = 4 - 3 = 1 > 0 \quad (\text{双曲型})
\]
特性方程式は
\[
\frac{dy}{dx} = \frac{B \pm \sqrt{D}}{A} = -2 \pm 1
\]
すなわち $\dfrac{dy}{dx} = -1$ または $\dfrac{dy}{dx} = -3$。それぞれ積分して
\[
y + x = c_1, \qquad y + 3x = c_2
\]
そこで
\[
\xi = y + x, \qquad \eta = y + 3x
\]
とおく。逆に $x = \dfrac{\eta-\xi}{2},\ y = \dfrac{3\xi-\eta}{2}$。連鎖律より
\[
u_x = u_\xi + 3u_\eta, \qquad u_y = u_\xi + u_\eta
\]
\[
u_{xx} = u_{\xi\xi} + 6u_{\xi\eta} + 9u_{\eta\eta}, \qquad
u_{xy} = u_{\xi\xi} + 4u_{\xi\eta} + 3u_{\eta\eta}, \qquad
u_{yy} = u_{\xi\xi} + 2u_{\xi\eta} + u_{\eta\eta}
\]
これを元の方程式に代入すると
\[
(u_{\xi\xi}+6u_{\xi\eta}+9u_{\eta\eta}) - 4(u_{\xi\xi}+4u_{\xi\eta}+3u_{\eta\eta}) + 3(u_{\xi\xi}+2u_{\xi\eta}+u_{\eta\eta})
\]
係数を整理すると、$u_{\xi\xi}$ の係数は $1-4+3=0$、$u_{\eta\eta}$ の係数は $9-12+3=0$、$u_{\xi\eta}$ の係数は $6-16+6=-4$ となるので
\[
-4u_{\xi\eta} = 0 \quad\Longrightarrow\quad u_{\xi\eta}=0
\]
これが標準形である。$u_{\xi\eta}=0$ を $\eta$ について積分すると $u_\xi = f'(\xi)$（$\eta$ に依らない関数）、さらに $\xi$ について積分して
\[
u = f(\xi) + g(\eta)
\]
（$f,g$ は任意の $C^2$ 関数）。$\xi,\eta$ をもとに戻すと
\[
\boxed{\; u(x,y) = f(y+x) + g(y+3x) \;}
\]
（検算: $u=f(y+x)$ とすると $u_{xx}=f'',\,u_{xy}=f'',\,u_{yy}=f''$ より $u_{xx}-4u_{xy}+3u_{yy}=(1-4+3)f''=0$。$u=g(y+3x)$ とすると $u_{xx}=9g'',\,u_{xy}=3g'',\,u_{yy}=g''$ より $(9-12+3)g''=0$。したがって重ね合わせにより一般解は方程式を満たす。）

\clearpage
\begin{mondaiboxS}{問3}
$u_{xx} - 2u_{xy} + u_{yy} = 0$ を標準形に変換し、初期条件 $u(x,0)=x^{2}$, $u_y(x,0)=0$ を満たす解を求めよ。
\end{mondaiboxS}
\kaitou
$A=1,\ B=-1,\ C=1$ であるから
\[
D = B^{2}-AC = 1-1 = 0 \quad (\text{放物型})
\]
特性方程式は
\[
\frac{dy}{dx} = \frac{B}{A} = -1
\]
（重根なので特性曲線は1本）。積分して $y+x=c$。そこで
\[
\xi = y+x, \qquad \eta = y
\]
（$\eta$ は $\xi,\eta$ が独立になるように任意に選んだもの）とおく。逆に $x=\xi-\eta,\ y=\eta$。連鎖律より
\[
u_x = u_\xi, \qquad u_y = u_\xi + u_\eta
\]
\[
u_{xx}=u_{\xi\xi}, \qquad u_{xy}=u_{\xi\xi}+u_{\xi\eta}, \qquad u_{yy}=u_{\xi\xi}+2u_{\xi\eta}+u_{\eta\eta}
\]
元の方程式に代入すると
\[
u_{\xi\xi} - 2(u_{\xi\xi}+u_{\xi\eta}) + (u_{\xi\xi}+2u_{\xi\eta}+u_{\eta\eta})
= (1-2+1)u_{\xi\xi} + (-2+2)u_{\xi\eta} + u_{\eta\eta} = u_{\eta\eta}
\]
よって標準形は
\[
u_{\eta\eta} = 0
\]
これを $\eta$ について2回積分すると
\[
u = f(\xi) + \eta\,g(\xi)
\]
（$f,g$ は任意関数）。$\xi=y+x,\ \eta=y$ に戻すと一般解は
\[
u(x,y) = f(x+y) + y\,g(x+y)
\]
次に初期条件を課す。$u_y = f'(x+y) + g(x+y) + y\,g'(x+y)$ であるから
\[
u(x,0) = f(x) = x^{2}
\]
\[
u_y(x,0) = f'(x) + g(x) = 0 \quad\Longrightarrow\quad g(x) = -f'(x) = -2x
\]
これらを一般解に代入して、$f(x+y)=(x+y)^2$, $g(x+y)=-2(x+y)$ より
\[
u(x,y) = (x+y)^{2} - 2y(x+y)
\]
展開すると $(x+y)^2 - 2y(x+y) = x^2+2xy+y^2-2xy-2y^2 = x^2 - y^2$ なので
\[
\boxed{\; u(x,y) = x^{2} - y^{2} \;}
\]
（検算: $u=x^2-y^2$ より $u_{xx}=2,\ u_{xy}=0,\ u_{yy}=-2$ で $u_{xx}-2u_{xy}+u_{yy}=2-0-2=0$ を満たす。また $u(x,0)=x^2$、$u_y=-2y$ より $u_y(x,0)=0$ もともに満たす。）

\end{document}
